% Irodov Part VI §6.1 problems 33--48 (end of §6.1).
% Source: Arihant PDF pp.245--246 (printed pp.236--237),
% Answers pp.349 (printed pp.340). Source scans establish
% #36=5.18e15 s^-1, #39=18.0 nm, #44=0.20 nm.
\atomicproblem{33}{236}
He$^+$ 离子从某一激发态连续发射两光子，
其波长分别为 $108.5$ 与 $30.4\,\mathrm{nm}$，
最终达到基态。求初始主量子数 $n$。
\exsolution
He$^+$ 基态结合能约 $E_b=54.4\,\mathrm{eV}$。
由跃迁能量
\[
 \Delta E=E_b\left(\frac1{n_f^2}-\frac1{n_i^2}\right)
 =\frac{hc}{\lambda}
\]
知 $30.4\,\mathrm{nm}$ 对应约
$40.8\,\mathrm{eV}$，
即 $2\to1$ 跃迁。
$108.5\,\mathrm{nm}$ 对应约
$11.43\,\mathrm{eV}$，因此前一跃迁满足
\[
 11.43=54.4\left(\frac14-\frac1{n^2}\right).
\]
解得 $1/n^2\simeq0.040$，
\[
 \boxed{n=5},\quad
 \text{跃迁顺序 }5\to2\to1.
\]
\excheck 两光子能量之和
$11.43+40.8\simeq52.2\,\mathrm{eV}$，
等于 $54.4(1-1/25)\,\mathrm{eV}$。

\atomicproblem{34}{236}
已知 He$^+$ 的巴耳末系和莱曼系第一条
（各自最长波长）谱线波长之差为
$\Delta\lambda=133.7\,\mathrm{nm}$，
求以\emph{角频率}记号定义的里德伯常数
$\mathcal R$。
\exsolution
类氢离子跃迁角频率写作
$\omega=\mathcal R Z^2(n_1^{-2}-n_2^{-2})$，
对 He$^+$ 取 $Z=2$。
第一条巴耳末为 $3\to2$，
第一条莱曼为 $2\to1$，因此
\[
 \lambda_B=\frac{2\pi c}{\mathcal RZ^2(1/4-1/9)}
 =\frac{72\pi c}{5\mathcal RZ^2},
\quad
 \lambda_L=\frac{2\pi c}{\mathcal RZ^2(1-1/4)}
 =\frac{8\pi c}{3\mathcal RZ^2}.
\]
相减得
\[
 \Delta\lambda=\frac{176\pi c}{15\mathcal RZ^2},
 \qquad
 \boxed{\mathcal R=
 \frac{176\pi c}{15Z^2\Delta\lambda}
 \simeq2.07\times10^{16}\,\mathrm{s^{-1}}}.
\]
\excheck 这里 $\mathcal R$ 是\emph{角频率常数}，
与常用波数里德伯常数
$R_\infty\simeq1.097\times10^7\,\mathrm{m^{-1}}$
相差因子 $2\pi c$。

\atomicproblem{35}{236}
某单电子类氢离子巴耳末系与莱曼系
第一谱线的波长差为 $59.3\,\mathrm{nm}$。
确定该离子。
\exsolution
由上一题可知，对于给定跃迁组合，
波长差按核电荷 $Z$ 缩放为
$\Delta\lambda_Z\propto1/Z^2$。
He$^+$（$Z=2$）时差为 $133.7\,\mathrm{nm}$，
于是
\[
 \left(\frac Z2\right)^2=
 \frac{133.7}{59.3},\quad
 Z\simeq3.00.
\]
因此为
\[
 \boxed{\mathrm{Li}^{2+}
 \quad\text{（双电离锂，核电荷数3）}}.
\]
\excheck 中性锂原子含3电子，不属于类氢单电子体系。

\atomicproblem{36}{236}
He$^+$ 的某一谱系，其最大与最小
跃迁角频率之差为
$\Delta\omega=5.18\times10^{15}\,\mathrm{s^{-1}}$。
求这个谱系的第一条（最长波长）谱线波长。
\exsolution
设谱系下能级为 $n$。
系限角频率为
$\omega_\infty=\mathcal R Z^2/n^2$，
第一线来自 $n+1\to n$，
$\omega_1=\mathcal RZ^2[n^{-2}-(n+1)^{-2}]$。
故间隔
\[
 \Delta\omega=\omega_\infty-\omega_1
 =\frac{\mathcal R Z^2}{(n+1)^2}.
\]
对 $Z=2$、$\mathcal R=2.07\times10^{16}\,\mathrm{s^{-1}}$，
得到 $(n+1)^2\simeq16$，所以 $n=3$。
第一线为 $4\to3$，
\[
 \boxed{\lambda_{43}=
 \frac{2\pi c}{\mathcal R\,4(1/9-1/16)}
 \simeq0.47\,\mu\mathrm m}.
\]
\excheck 题中是两端\emph{角频率}之差，
不可用普通频率常数遗漏 $2\pi$。

\atomicproblem{37}{236}
某类氢离子巴耳末系第三条谱线波长
为 $108.5\,\mathrm{nm}$。
求该离子基态电子结合能。
\exsolution
巴耳末第三线为 $5\to2$，
对应光子能量
\[
 \frac{hc}{\lambda}
 =E_b\left(\frac14-\frac1{25}\right)
 =\frac{21}{100}E_b.
\]
所以
\[
 \boxed{E_b=\frac{100}{21}
 \frac{hc}{108.5\,\mathrm{nm}}
 \simeq54.4\,\mathrm{eV}}.
\]
该值对应 He$^+$。
\excheck 不同谱线序数须区别于主量子数；
第三条巴耳末线的上能级为5。

\atomicproblem{38}{236}
中性氦原子基态第一个电子的结合能
（第一电离能）为
$E_0=24.6\,\mathrm{eV}$。
求移去氦原子全部两个电子所需总能量。
\exsolution
先移去第一个电子，消耗
$24.6\,\mathrm{eV}$；
剩余 He$^+$ 是核电荷为2的类氢离子，
其基态第二电离能为
$4\times13.6=54.4\,\mathrm{eV}$。
所以
\[
 \boxed{E_{\rm total}=24.6+54.4
 =79.0\,\mathrm{eV}}.
\]
\excheck 不能把第一电离能简单乘2，
因为第一次电离前后电子间排斥的影响不同。

\atomicproblem{39}{236}
波长 $\lambda=18.0\,\mathrm{nm}$ 的辐射
使起初静止的 He$^+$ 离子基态发生光电离。
求光电子最大速度。
\exsolution
单个入射光子能量
$hc/\lambda\simeq68.88\,\mathrm{eV}$，
He$^+$ 的结合能 $E_b=54.4\,\mathrm{eV}$。
忽略重核反冲，
电子动能
$K=hc/\lambda-E_b\simeq14.48\,\mathrm{eV}$。
故
\[
 \boxed{v_{\max}=
 \sqrt{\frac{2(hc/\lambda-E_b)}{m_e}}
 \simeq2.26\times10^6\,\mathrm{m/s}}.
\]
\excheck 原扫描题面为 \textbf{$18.0\,\mathrm{nm}$}，
不是 PDF 抽取文本易误识别的 $180\,\mathrm{nm}$；
后者光子能量不足以电离基态 He$^+$。

\atomicproblem{40}{236}
基态氢原子 A 以初始动能 $T$ 正面撞击
另一个起初静止的基态氢原子 B。
求使其中一个原子至少激发到
$n=2$、进而能辐射光子的最小 $T$。
\exsolution
使一个氢原子从基态激发到 $n=2$，
需要内部激发能
$\Delta=13.6(1-1/4)=10.2\,\mathrm{eV}$。
两个原子质量相等 $M$，
实验室系中可转化为内部激发的质心相对运动能量
为
\[
 T_{\rm cm}=\frac12\mu v^2
 =\frac12\frac M2\,v^2
 =\frac12T,\qquad T=\frac12Mv^2.
\]
所以临界条件 $T/2=\Delta$，
\[
 \boxed{T_{\min}=2\Delta\simeq20.4\,\mathrm{eV}}.
\]
\excheck 不能仅取 $10.2\,\mathrm{eV}$，
因为碰后质心必须继续运动，
入射动能不可能全部转为内部激发。

\atomicproblem{41}{236}
起初静止的氢原子放出第一条莱曼线光子。
求氢原子获得的反冲速度。
\exsolution
第一莱曼线来自 $2\to1$，
光子能量
$\varepsilon\simeq10.2\,\mathrm{eV}$，
动量 $p_\gamma=\varepsilon/c$。
若原子起初静止，由总动量守恒，
原子获得等大反向动量，故
\[
 \boxed{v_{\mathrm H}=
 \frac{\varepsilon}{M_{\mathrm H}c}
 \simeq3.25\,\mathrm{m/s}}.
\]
\excheck 原子质量使用氢\emph{原子}质量；
光子的质量不是 $E/c^2$ 再去用非相对论动量公式求速度。

\atomicproblem{42}{236}
承上题，求实际发出光子的能量
与不计反冲时的 $2\to1$ 内能差之间的相对差，
以百分数表示。
\exsolution
跃迁内能差记为 $\Delta E$。
发射后原子后退动能
$K_{\rm recoil}=p_\gamma^2/(2M)
=E_\gamma^2/(2Mc^2)$。
能量守恒
$\Delta E=E_\gamma+E_\gamma^2/(2Mc^2)$；
因 $\Delta E\ll Mc^2$，
一阶近似
\[
 \frac{\Delta E-E_\gamma}{\Delta E}
 \simeq\frac{\Delta E}{2Mc^2}.
\]
以 $\Delta E=10.2\,\mathrm{eV}$、
$Mc^2\simeq938.8\,\mathrm{MeV}$，
\[
 \boxed{\text{相对差}\simeq
 5.43\times10^{-9}
 =5.43\times10^{-7}\%\approx0.55\times10^{-6}\%}.
\]
\excheck 这是一个极小的\emph{红移}校正；
将比值改写百分数时必须乘100。

\atomicproblem{43}{236}
起初静止的 He$^+$ 发出第一莱曼线光子，
该光子随后电离起初静止的基态氢原子。
求被释放光电子的速度，忽略重核反冲。
\exsolution
He$^+$ 的 $2\to1$ 光子能量为
$4(10.2)=40.8\,\mathrm{eV}$。
氢基态电离能为
$13.6\,\mathrm{eV}$。
由能量守恒，
光电子获得动能
$T_e=40.8-13.6=27.2\,\mathrm{eV}$，
因而
\[
 \boxed{v_e=\sqrt{\frac{2T_e}{m_e}}
 \simeq3.09\times10^6\,\mathrm{m/s}}.
\]
\excheck 光子具有离散跃迁能，
不能将 He$^+$ 的 $40.8\,\mathrm{eV}$ 误作光电子动能；
必须扣除氢的电离能。

\atomicproblem{44}{236}
沿与原子速度成 $\theta=45^\circ$
的方向观察激发氢原子的第一莱曼线，
发现波长位移 $\Delta\lambda=0.20\,\mathrm{nm}$。
估算氢原子速率。
\exsolution
在 $v\ll c$ 时，相对论多普勒公式的一阶形式为
\[
 \frac{|\Delta\lambda|}{\lambda_0}
 \simeq\frac{v|\cos\theta|}{c}.
\]
第一莱曼线
$\lambda_0\simeq121.6\,\mathrm{nm}$，
因此
\[
 \boxed{v\simeq
 \frac{c\,|\Delta\lambda|}
 {\lambda_0|\cos45^\circ|}
 \simeq6.98\times10^5\,\mathrm{m/s}}.
\]
\excheck 位移的符号由接近或远离决定；
题给只问速度大小，故取绝对值。

\atomicproblem{45}{236--237}
应用玻尔--索末菲作用量量子化
$\oint p\,dq=2\pi n\hbar$，
求粒子在：(a) 长度 $\ell$ 的无限深一维方阱；
(b) 半径 $r$ 的圆周；
(c) $U(x)=\alpha x^2/2$ 的一维谐振势；
(d) $U(r)=-\alpha/r$ 的中心引力型势作圆周运动时，
允许的能量。
\exsolution
(a) 粒子在方阱两侧来回反射，一个周期走过 $2\ell$，
故 $2p\ell=2\pi n\hbar$，
\[
 \boxed{E_n=\frac{n^2\pi^2\hbar^2}{2m\ell^2}}.
\]
(b) 沿圆周 $\oint p\,dq=2\pi rp=2\pi L$，
所以 $L=n\hbar$，有
\[
 \boxed{E_n=\frac{n^2\hbar^2}{2mr^2}}.
\]
(c) 振子频率 $\Omega=\sqrt{\alpha/m}$，
经典一周作用量
$\oint p\,dx=2\pi E/\Omega$，于是
\[
 \boxed{E_n=n\hbar\Omega=n\hbar\sqrt{\alpha/m}}.
\]
(d) 圆周条件 $mv^2/r=\alpha/r^2$，
加上 $L=mvr=n\hbar$，
解出
$r=n^2\hbar^2/(m\alpha)$，
能量为
\[
 \boxed{E_n=-\frac{m\alpha^2}{2n^2\hbar^2}}.
\]
\excheck 这是\emph{早期作用量量子化}推导；
(c) 缺少精确量子谐振子的零点能
$\hbar\Omega/2$，不应把这两种理论混为一谈。

\atomicproblem{46}{237}
考虑氢原子质子的运动，
求基态电子结合能及里德伯常数的
约化质量修正，并估计忽略核运动的相对误差。
\exsolution
电子与质子的相对运动等价于约化质量
\[
 \mu=\frac{m_em_p}{m_e+m_p}
\]
的单粒子库仑问题。
因玻尔结合能与里德伯常数都正比于质量，
\[
 \boxed{E_b=\frac{\mu(k_ee^2)^2}{2\hbar^2}},\qquad
 \boxed{\mathcal R=
 \frac{\mu(k_ee^2)^2}{2\hbar^3}},
\]
其中 $\mathcal R$ 是角频率形式，
若使用波数里德伯常数需除以 $2\pi c$。
忽略质子运动等于用 $m_e$ 代替 $\mu$，
相对高估量为
\[
 \frac{m_e-\mu}{\mu}=\frac{m_e}{m_p}
 \simeq5.45\times10^{-4}.
\]
故
\[
 \boxed{\text{相对误差}\simeq0.0545\%}.
\]
\excheck 对真实氢原子这一修正很小，
但分析同位素光谱线位移时不可遗漏。

\atomicproblem{47}{237}
对普通氢 H 和重氢 D，
计算：(a) 基态电子结合能之差；
(b) 第一莱曼线波长之差。
\exsolution
两体系电子质量相同，核质量分别为
$m_p$、$m_d$，约化质量
$\mu_A=m_em_A/(m_e+m_A)$。
结合能及跃迁频率都与 $\mu$ 成正比，
故
\[
 \boxed{E_{b,D}-E_{b,H}
 =13.6057\,\mathrm{eV}\,
 \frac{\mu_D-\mu_H}{m_e}
 \simeq3.7\,\mathrm{meV}}.
\]
第一莱曼线波长
$\lambda_{21,A}\propto1/\mu_A$，
因此
\[
 \boxed{\lambda_{21,H}-\lambda_{21,D}
 \simeq33\,\mathrm{pm}}.
\]
\excheck 重氢核较重，约化质量更接近 $m_e$，
因此 D 的结合能稍大、莱曼波长稍短。

\atomicproblem{48}{237}
估计以下两个类氢系统基态时
两粒子间距、结合能和第一莱曼线波长：
(a) 质子与带电量 $-e$、质量 $207m_e$ 的负介子
组成的介子氢；
(b) 电子与正电子组成的正电子素。
\exsolution
对电荷模为 $e$ 的两粒子库仑体系，
只需把电子质量换成约化质量
$\mu=m_1m_2/(m_1+m_2)$。
基态类玻尔半径、结合能、第一莱曼线依次为
\[
 a=\frac{4\pi\varepsilon_0\hbar^2}{\mu e^2},
 \quad E_b=\frac{\mu(k_ee^2)^2}{2\hbar^2},
 \quad\lambda_{21}=\frac{hc}{3E_b/4}.
\]
(a) 对 $\mu_{\mu p}=
(207m_e)m_p/(207m_e+m_p)\simeq186m_e$，
\[
 \boxed{a\simeq0.285\,\mathrm{pm},\quad
 E_b\simeq2.53\,\mathrm{keV},\quad
 \lambda_{21}\simeq0.65\,\mathrm{nm}}.
\]
(b) 正电子素约化质量恰为 $\mu=m_e/2$，
所以
\[
 \boxed{a\simeq106\,\mathrm{pm},\quad
 E_b\simeq6.80\,\mathrm{eV},\quad
 \lambda_{21}\simeq243\,\mathrm{nm}}.
\]
\excheck 第(a)问不能直接用介子质量 $207m_e$
代替约化质量，质子质量有限会产生明显修正。
